\documentclass[11pt]{article}
\usepackage[letterpaper,margin=1in]{geometry}
\usepackage{amsmath,amssymb,amsthm,mathtools}
\usepackage{booktabs,array,longtable}
\usepackage{xcolor}
\usepackage{microtype}
\usepackage[hidelinks]{hyperref}
\hypersetup{
  pdftitle={Explicit nonlinear SU(2) geometry on a four-dimensional lattice block},
  pdfsubject={Quaternion, incidence, local gauge-slice, and action-norm certificates}
}

\newtheorem{theorem}{Theorem}[section]
\newtheorem{proposition}[theorem]{Proposition}
\newtheorem{lemma}[theorem]{Lemma}
\newtheorem{corollary}[theorem]{Corollary}
\theoremstyle{definition}
\newtheorem{definition}[theorem]{Definition}
\newtheorem{remark}[theorem]{Remark}

\newcommand{\SU}{\mathrm{SU}(2)}
\newcommand{\su}{\mathfrak{su}(2)}
\newcommand{\R}{\mathbb{R}}
\newcommand{\norm}[1]{\left\lVert #1\right\rVert}
\newcommand{\ip}[2]{\left\langle #1,#2\right\rangle}
\newcommand{\dd}{\mathrm{d}}
\newcommand{\statusP}{\textcolor{green!45!black}{\textsf{P}}}
\newcommand{\statusN}{\textcolor{blue!65!black}{\textsf{N}}}
\newcommand{\statusO}{\textcolor{red!70!black}{\textsf{O}}}

\title{Explicit Nonlinear $\SU$ Geometry on a Four-Dimensional Lattice Block\\
\large Quaternion bounds, incidence geometry, a local gauge slice, and a five-norm initializer}
\author{Mathematical research dossier}
\date{20 July 2026}

\begin{document}
\maketitle

\begin{abstract}
We give explicit finite-dimensional constants for the nonlinear plaquette-log
map on the side-two four-dimensional cubical block.  The argument represents
$\SU$ by unit quaternions, bounds the first three derivatives of a product
logarithm by an absolutely convergent Banach-algebra series, and combines
these bounds with the exact cubical incidence constants.  On the logarithmic
Coulomb slice the nonlinear curvature map is quantitatively bi-Lipschitz.  A
separate contraction theorem constructs the local Coulomb representative,
modulo residual constant conjugation.  Composing these geometric estimates
with the exact heat-kernel derivative certificate produces all five local
action-norm components required by a one-step retained/fiber marginalization.
For $t\leq 0.01$, four-link radius $0.003$, and centered fiber radius $0.001$,
the scaled Hessian remainder is at most $0.747472$, the centered-fiber
transport coefficient is at most $0.266017$, and the generated reference-base
margin is at least $0.403230$.

The same certificate isolates the next mathematical obstruction.  At the
upper heat time, the one-plaquette probability outside the certified
curvature radius $0.006$ exceeds $0.9999$.  Thus the finite local geometry is
closed, while a probabilistically useful chart and a dependent bad-block
polymer estimate remain necessary for a full regulated RG theorem.
\end{abstract}

\tableofcontents

\section{Result and mathematical role}

The exact $\SU$ heat-kernel action gives a gauge-invariant plaquette density,
as developed in the lattice formulation of Menotti and Onofri
\cite{MenottiOnofri1981}.  A constructive block transformation requires more
than the radial one-plaquette formula: the action must be transported through
noncommuting link coordinates, gauge directions must be removed, and local
derivative estimates must be aggregated over the block.  These are the
finite-dimensional operations treated here.  Their organization follows the
small-field side of gauge-covariant block RG constructions such as
Balaban's \cite{Balaban1984,Balaban1987}.

Let $B$ be the cubical complex with vertices $\{0,1,2\}^4$.  Write $B_1$ for
its positively oriented links and $B_2$ for its positively oriented
plaquettes.  A link-log field is
\[
 A=(A_e)_{e\in B_1}\in \su^{B_1}\cong\R^{648},
 \qquad U_e=\exp(A_e).
\]
For $p=(x;\mu,\nu)$, $\mu<\nu$, define
\begin{equation}\label{eq:holonomy}
 U_p(A)=U_{x,\mu}U_{x+\hat\mu,\nu}
 U_{x+\hat\nu,\mu}^{-1}U_{x,\nu}^{-1},
 \qquad X_p(A)=\log U_p(A).
\end{equation}
The scalar incidence matrices are $d_0:C^0(B)\to C^1(B)$ and
$d_1:C^1(B)\to C^2(B)$.  They act diagonally on the three color components.

\begin{theorem}[Calibrated nonlinear block certificate]\label{thm:main}
Use heat generator $\frac12\Delta_{\SU}$ and let
\[
 t_{\max}=10^{-2},\qquad r=3\cdot10^{-3},\qquad
 \rho=2r=6\cdot10^{-3},\qquad r_F=10^{-3}.
\]
On the convex Coulomb chart
\[
 \mathcal C_r=\left\{A\in\ker d_0^T:
   \norm{A|_{\partial p}}_2\le r\text{ for every }p\in B_2\right\},
\]
the following hold.
\begin{enumerate}
\item The map $X=(X_p)_{p\in B_2}$ is injective and satisfies
\[
 \norm{X(A)-X(B)}_2\ge 1.354892\,\norm{A-B}_2,
 \qquad A,B\in\mathcal C_r.
\]
Consequently its inverse on its image is $0.738066$-Lipschitz.
\item Mean-zero vertex gauge parameters of norm at most $6\cdot10^{-4}$
give a contraction with coefficient at most $0.051367$.  Every input link-log
field of norm at most $2\cdot10^{-4}$ has a unique Coulomb representative in
that parameter ball, modulo constant conjugation, and the representative has
norm at most $0.0026<r$.
\item For the pure plaquette action
\[
 S_t(A)=\sum_{p\in B_2}-\log p_t\bigl(U_p(A)\bigr),
 \qquad 0<t\le t_{\max},
\]
the global scaled Hessian remainder relative to
$\norm{d_1A}^2/(2t)$ is at most
\[
 \eta=0.747472,
\]
and the global scaled third derivative is at most
$\tau=295.100$.
\item In Hodge-whitened retained coordinates and an orthonormal conditional
fiber coordinate, the five good-chart budgets may be chosen as
\begin{align*}
 \varepsilon_F&=0.747472,&
 \varepsilon_{FY}&=0.528543,&
 \varepsilon_Y&=0.373736,\\
 \tau_F&=295.100,&
 \tau_{FY}&=208.668.
\end{align*}
With reference curvatures $\lambda_F=2$, $\lambda_Y=1$, these give
\[
 K_F\le0.266017<1,
 \qquad
 \lambda_Y-\eta_{\mathrm{good}}\ge0.403230>0.
\]
\end{enumerate}
These conclusions concern the declared local chart and its pure plaquette
action.  The density contributions of the link Haar measure and the gauge
slice, together with the chart complement, remain separate analytic inputs.
\end{theorem}

The proof occupies Sections~\ref{sec:quaternion}--\ref{sec:action}.

\section{Quaternion product-log calculus}\label{sec:quaternion}

Identify $\SU$ with the unit quaternions.  For $a\in\R^3$ regarded as a pure
imaginary quaternion,
\[
 \exp(a)=\left(\cos\norm a,\frac{\sin\norm a}{\norm a}a\right).
\]
Quaternion norm is multiplicative.  The principal logarithm of a unit
quaternion $(q_0,q)$ before the antipode is
\[
 \log(q_0,q)=\frac{\operatorname{atan2}(\norm q,q_0)}{\norm q}\,q.
\]

\begin{lemma}[Banach logarithm derivatives]\label{lem:log}
In a unital normed algebra with submultiplicative norm, let $z=1+h$ with
$\norm h\le\delta<1$.  Then
\begin{equation}\label{eq:logbounds}
 \norm{D\log z}\le\frac1{1-\delta},\qquad
 \norm{D^2\log z}\le\frac1{(1-\delta)^2},\qquad
 \norm{D^3\log z}\le\frac2{(1-\delta)^3}.
\end{equation}
\end{lemma}

\begin{proof}
The series
\[
 \log(1+h)=\sum_{n\ge1}\frac{(-1)^{n+1}}n h^n
\]
converges absolutely.  The first derivative of $h^n$ contains $n$ ordered
insertions, the second contains $n(n-1)$, and the third contains
$n(n-1)(n-2)$.  After division by $n$, summing the norm majorants gives
\[
 \sum_{n\ge1}\delta^{n-1}=\frac1{1-\delta},\quad
 \sum_{n\ge2}(n-1)\delta^{n-2}=\frac1{(1-\delta)^2},
\]
and
\[
 \sum_{n\ge3}(n-1)(n-2)\delta^{n-3}=\frac2{(1-\delta)^3}.
\]
\end{proof}

\begin{lemma}[Product derivatives]\label{lem:product}
Let
\[
 P(a_1,\ldots,a_m)=\exp(a_1)\cdots\exp(a_m)
\]
with the direct-sum Euclidean norm on the inputs.  For $k=1,2,3$,
\begin{equation}\label{eq:productderivatives}
 \norm{D^kP}\le m^{k/2}.
\end{equation}
Moreover, on the input ball of radius $r$,
$\norm{P-1}\le\sqrt m\,r$.
\end{lemma}

\begin{proof}
The iterated Duhamel formula for $D^k\exp(a)$ is a sum over permutations of
simplex integrals.  Every exponential factor has quaternion norm one and the
sum of the simplex volumes is one, so
$\norm{D^k\exp(a)[u_1,\ldots,u_k]}\le\prod_j\norm{u_j}$.
Applying the product rule and then Cauchy--Schwarz to each differentiated
input gives $m^{k/2}$.  The final estimate follows by telescoping the product
from the identity and using
$\norm{\exp(a_i)-1}\le\norm{a_i}$.
\end{proof}

Combining Lemmas~\ref{lem:log} and \ref{lem:product} gives a closed formula.
For $\delta=\sqrt m,r<1$, set
\begin{align}
 M_1^{(m)}(r)&=\frac{\sqrt m}{1-\delta},\label{eq:M1}\\
 M_2^{(m)}(r)&=m\left(\frac1{(1-\delta)^2}+\frac1{1-\delta}\right),\label{eq:M2}\\
 M_3^{(m)}(r)&=\frac{2m^{3/2}}{(1-\delta)^3}
 +\frac{3m^{3/2}}{(1-\delta)^2}
 +\frac{m^{3/2}}{1-\delta}.
\label{eq:M3}
\end{align}

\begin{proposition}[Four-factor plaquette chart]\label{prop:plaquette}
For the oriented plaquette map \eqref{eq:holonomy}, if the joint norm of its
four link inputs is at most $r$ and $2r<1$, then
\[
 \norm{DX_p}\le M_1,\qquad
 \norm{D^2X_p}\le M_2,\qquad
 \norm{D^3X_p}\le M_3,
\]
where
\begin{align*}
 M_1&=\frac2{1-2r},\\
 M_2&=\frac4{(1-2r)^2}+\frac4{1-2r},\\
 M_3&=\frac{16}{(1-2r)^3}
       +\frac{24}{(1-2r)^2}+\frac8{1-2r}.
\end{align*}
If $L_p=d_1|_p$ is the oriented linear boundary sum, then
\begin{equation}\label{eq:localdeviation}
 \norm{DX_p(A)-L_p}\le M_2r,\qquad \norm{X_p(A)}\le2r.
\end{equation}
\end{proposition}

\begin{proof}
Signs in the two inverse links are absorbed into $a_i\mapsto-a_i$, an
isometry.  Equations \eqref{eq:M1}--\eqref{eq:M3} with $m=4$ give the first
three claims.  At the origin $DX_p(0)=L_p$; integrating $D^2X_p$ along the
radial segment gives the derivative deviation.  The last inequality is the
triangle inequality for the bi-invariant geodesic distance on $\SU$.
\end{proof}

\section{Exact block incidence and quotient geometry}\label{sec:incidence}

Direct cubical enumeration gives
\[
 |B_0|=81,\qquad |B_1|=216,\qquad |B_2|=216,\qquad d_1d_0=0.
\]
Every plaquette has four boundary links.  A link lies in at most six
plaquettes, and a vertex meets at most eight links.  Therefore, if $E_p$ is
the extraction of the four link variables around $p$,
\begin{equation}\label{eq:incidence}
 \sum_{p\in B_2}\norm{E_ph}^2\le6\norm h^2.
\end{equation}

The interval with two edges has vertex-Laplacian spectrum $\{0,1,3\}$.
The Cartesian product rule gives the four-dimensional spectra by summing
four interval eigenvalues.  Consequently
\begin{equation}\label{eq:hodge}
 \norm{d_0}=\norm{d_1}=\sqrt{12},\qquad
 \norm{d_1h}\ge\sqrt2\norm h\quad(h\in\ker d_0^T).
\end{equation}
The ranks are
\[
 \operatorname{rank}d_0=80,\qquad
 \operatorname{rank}d_1=136.
\]
The 24 displayed boundary-face fluxes have rank 17 on the physical quotient,
leaving scalar fiber dimension $136-17=119$.  Tensoring with $\su$ gives
retained and fiber dimensions 51 and 357.

\begin{theorem}[Nonlinear curvature coordinates]\label{thm:bilipschitz}
On the plaquettewise radius-$r$ Coulomb chart, define
\[
 \kappa(r)=\sqrt6,M_2(r)r.
\]
Then
\begin{equation}\label{eq:globaldeviation}
 \norm{DX(A)-d_1}\le\kappa(r).
\end{equation}
If $\kappa(r)<\sqrt2$, the curvature map is injective on the chart and
\begin{equation}\label{eq:bilipschitz}
 \norm{X(A)-X(B)}\ge(\sqrt2-\kappa(r))\norm{A-B}.
\end{equation}
\end{theorem}

\begin{proof}
Apply \eqref{eq:localdeviation} to every plaquette and sum squares.  The
incidence estimate \eqref{eq:incidence} yields
\eqref{eq:globaldeviation}.  For $A,B$ in the convex chart,
\[
 X(A)-X(B)=d_1(A-B)+
 \int_0^1\bigl(DX(B+s(A-B))-d_1\bigr)(A-B)\,\dd s.
\]
The Hodge lower bound \eqref{eq:hodge} and the triangle inequality give
\eqref{eq:bilipschitz}.
\end{proof}

For $r=0.003$,
\[
 M_1=2.012073,\quad M_2=8.072581,\quad M_3=48.630394,\quad
 \kappa=0.059322,
\]
so $\sqrt2-\kappa=1.354892$.

\section{A quantitative local Coulomb slice}\label{sec:slice}

Let
\[
 V_0=\left\{\phi\in(\R^3)^{B_0}:\sum_{x\in B_0}\phi_x=0\right\}
\]
be the nonconstant gauge-parameter space.  Its orthogonal complement is the
residual constant conjugation.  For a link $e=(x,y)$ set
\begin{equation}\label{eq:gaugeaction}
 A_e^\phi=\log\left(\exp(\phi_x)\exp(A_e)\exp(-\phi_y)\right)
\end{equation}
and define $F_A(\phi)=d_0^TA^\phi\in V_0$.  At the origin,
\[
 D_\phi F_0(0)=-\Delta_0,\qquad \Delta_0=d_0^Td_0|_{V_0}.
\]
The product spectrum gives $\norm{\Delta_0^{-1}}=1$.

\begin{theorem}[Local gauge contraction]\label{thm:gauge}
Suppose $\norm A\le a$ and seek $\phi\in V_0$ with $\norm\phi\le s$.  Put
\[
 R=\sqrt{a^2+2s^2},\qquad
 q=\sqrt{96}\,M_2^{(3)}(R)R.
\]
If $\sqrt3R<1$, $q<1$, and
\begin{equation}\label{eq:ballmap}
 a\le(1-q)s,
\end{equation}
then $F_A(\phi)=0$ has a unique solution in the radius-$s$ ball of $V_0$.
The transformed link field satisfies
\begin{equation}\label{eq:fieldoutput}
 \norm{A^\phi}\le a+4s.
\end{equation}
\end{theorem}

\begin{proof}
Consider
\[
 T_A(\phi)=\phi+\Delta_0^{-1}F_A(\phi).
\]
At zero, $A^0=A$, so
$\norm{T_A(0)}=\norm{\Delta_0^{-1}d_0^TA}\le a$.
The three-factor product-log bound controls the derivative remainder on every
edge by $M_2^{(3)}(R)R$.  An increment in a vertex field is extracted at the
two endpoints of each edge.  Since the maximum vertex degree is eight, the
global extraction norm is at most $\sqrt8$.  Composing with
$\norm{d_0^T}=\sqrt{12}$ and $\norm{\Delta_0^{-1}}=1$ gives
$\norm{DT_A}\le q$.  Condition \eqref{eq:ballmap} makes the radius-$s$ ball
invariant.  Banach's fixed-point theorem gives existence and uniqueness.
Finally, the geodesic triangle inequality on each edge and
\[
 \sum_{(x,y)\in B_1}(\norm{\phi_x}+\norm{\phi_y})^2
 \le16\norm\phi^2
\]
give \eqref{eq:fieldoutput}.
\end{proof}

For $a=0.0002$ and $s=0.0006$, the theorem gives
\[
 R=0.000871780,\quad q=0.0513663,\quad
 (1-q)s-a=0.000369180,\quad a+4s=0.0026.
\]
Thus the constructed Coulomb representative enters the radius-$0.003$
action chart.

\section{Heat-kernel action and the five-norm initializer}\label{sec:action}

Write the local heat-kernel action as
\[
 s_t(x)=\frac{\norm x^2}{2t}+R_t(x)+c(t).
\]
On $\norm x\le\rho$, Stage 07 supplies
\begin{equation}\label{eq:radialbudgets}
 t\norm{DR_t}\le g,\qquad
 t\norm{D^2R_t}\le\varepsilon,\qquad
 t\norm{D^3R_t}\le\tau_{\mathrm{HK}}.
\end{equation}
For $t\le0.01$ and $\rho=0.006$, the explicit values are
\[
 g\le2.00002\cdot10^{-5},\qquad
 \varepsilon\le0.00333336,\qquad
 \tau_{\mathrm{HK}}\le3.20004\cdot10^{-5}.
\]

\begin{proposition}[Refined block aggregation]\label{prop:action}
Let $\kappa$, $M_1$, $M_2$, and $M_3$ be the chart constants above.  The
scaled Hessian remainder of $S_t$ relative to
$\norm{d_1A}^2/(2t)$ is bounded by
\begin{equation}\label{eq:eta}
 \eta=
 \kappa(2\sqrt{12}+\kappa)
 +\varepsilon(\sqrt{12}+\kappa)^2
 +6(\rho+g)M_2.
\end{equation}
The scaled third derivative is bounded by
\begin{equation}\label{eq:tau}
 \tau=6\left[
 \tau_{\mathrm{HK}}M_1^3
 +3(1+\varepsilon)M_1M_2
 +(\rho+g)M_3
 \right].
\end{equation}
\end{proposition}

\begin{proof}
For the principal quadratic part,
\[
 D^2\frac12\norm X^2=DX^TDX+\ip{X}{D^2X}.
\]
The first difference from $d_1^Td_1$ is at most
$\kappa(2\sqrt{12}+\kappa)$.  The second is at most $6\rho M_2$ by the
incidence estimate.  The Hessian term from $R_t\circ X$ is bounded by
$\varepsilon\norm{DX}^2\le\varepsilon(\sqrt{12}+\kappa)^2$, and its gradient
term contributes at most $6gM_2$.  This proves \eqref{eq:eta}.

The third-order chain rule on one plaquette gives
\[
 \tau_{\mathrm{HK}}M_1^3+3(1+\varepsilon)M_1M_2
 +(\rho+g)M_3.
\]
For unit global directions $h,k,\ell$,
\[
 \sum_p\norm{E_ph}\norm{E_pk}\norm{E_p\ell}\le6,
\]
which proves \eqref{eq:tau}.
\end{proof}

Substitution yields
\[
 \eta=0.747471288,\qquad \tau=295.099754.
\]

\subsection{Retained/fiber coordinates}

Let $H=d_1^Td_1$ on the physical scalar quotient.  Let $C$ be the rank-17
retained coarse-flux map and $N$ an orthonormal basis of $\ker C$.  The
conditional-mean subspace $H^{-1}\operatorname{ran}C^T$ is $H$-orthogonal to
$N$.  Choose $M$ on that subspace so that
\begin{equation}\label{eq:whiten}
 M^THM=I,\qquad M^THN=0,\qquad N^THN\succeq2I.
\end{equation}
Since $H\succeq2I$ on the quotient,
\begin{equation}\label{eq:Mbound}
 \norm M\le\frac1{\sqrt2},\qquad \norm N=1.
\end{equation}
The same statements hold after tensoring with the three color components.

Writing $A=My+Nz$, the global Hessian and third-order bounds therefore give
\begin{equation}\label{eq:five}
 \varepsilon_F=\eta,\quad
 \varepsilon_{FY}=\frac\eta{\sqrt2},\quad
 \varepsilon_Y=\frac\eta2,\quad
 \tau_F=\tau,\quad
 \tau_{FY}=\frac\tau{\sqrt2}.
\end{equation}
With $r_F=0.001$, define
\begin{align*}
 \mu&=2-\varepsilon_F,\\
 h&=\varepsilon_{FY}/\mu,\\
 \chi&=r_F(\tau_{FY}+\tau_Fh),\\
 K_F&=\chi/\mu,\\
 \eta_{\mathrm{good}}&=\varepsilon_Y+\varepsilon_{FY}^2/\mu.
\end{align*}
The calibrated values are
\[
 \mu=1.252529,\quad h=0.421980,\quad
 K_F=0.266017,\quad \eta_{\mathrm{good}}=0.596770.
\]
This proves the action-norm part of Theorem~\ref{thm:main}.

\section{Probability overlap and the next inequality}

The chart is analytically nonempty, but an RG proof also needs its complement
to have a summable activity.  The one-plaquette exceptional probability can
be computed from the exact character expansion.  At the calibrated endpoint,
\begin{equation}\label{eq:probability}
 \Pr_{t=0.01}\{\theta(U_p)\ge0.006\}
 \in[0.9999423268758832,\,0.9999423268759137].
\end{equation}
The omitted-character bound after character 200 is below $1.2\cdot10^{-86}$;
a separate floating allowance is included in the displayed interval.

Equation \eqref{eq:probability} prevents a full-measure interpretation of the
local certificate.  The quantitative next target is an overlap window in
parameters $(t,r)$ satisfying simultaneously
\begin{equation}\label{eq:overlap}
 \eta(t,r)<2,\qquad K_F(t,r,r_F)<1,\qquad
 \norm{\mathfrak z_{\mathrm{bad}}(t,r)}_{\mathrm{poly}}<1,
\end{equation}
where the last expression is a dependent, gauge-invariant block-polymer norm,
not a product of one-plaquette probabilities.  An $\SU$-specific covariant
derivative calculus can improve the first two inequalities by exploiting
compact-group cancellations absent from the Banach-log majorant.  The third
requires shared-link dependence, Bianchi constraints, boundary holonomy,
center flux, and topological sectors.

\section{Verification and epistemic boundary}

The accompanying implementation performs the following independent checks:
\begin{itemize}
\item exact cubical enumeration, ranks, incidences, and numerical spectral
cross-checks;
\item quaternion exponential/logarithm round trips and commuting plaquette
identities;
\item random noncommutative plaquette samples against the analytic Taylor
remainder;
\item gauge invariance of the plaquette class angle and numerical convergence
of the Coulomb fixed-point map;
\item cross-check of the five local budgets against the separate Stage 06
one-step calculator;
\item hostile rejection of cut-locus, excessive-radius, lost-margin, and
false mass-gap promotions.
\end{itemize}

The theorem proves a local finite-block statement, classified at claim level
C2.  The following bridges remain open:
\begin{longtable}{>{\raggedright\arraybackslash}p{0.52\textwidth}c>{\raggedright\arraybackslash}p{0.32\textwidth}}
\toprule
Claim & Status & Missing object\\
\midrule
Quaternion, incidence, local inverse, and local Coulomb certificates & \statusP & Proved in this note\\
Five-component pure-plaquette good-chart norm & \statusP & Proved in this note\\
Executable finite-block constants and hostile tests & \statusN & Reproducible code and JSON\\
Link-Haar and gauge-slice determinant in the same norm & \statusO & Density derivative estimates\\
Dependent chart-complement polymer activity & \statusO & Probability-overlap theorem\\
Cutoff- and volume-uniform repeated RG & \statusO & Quasilocal interaction recursion\\
Reflection-positive OS continuum Yang--Mills theory & \statusO & Construction and theory identity\\
Positive Hamiltonian mass gap & \statusO & Basin entry and spectral transfer\\
\bottomrule
\end{longtable}

No finite-block coercivity constant is a physical mass.  The result advances
the constructive program by making the first nonlinear representation bridge
explicit and by locating its first quantitative failure at the interface
between local chart size and exceptional-set probability.

\begin{thebibliography}{9}
\bibitem{MenottiOnofri1981}
P.~Menotti and E.~Onofri,
\emph{The action of $SU(N)$ lattice gauge theory in terms of the heat kernel
on the group manifold}, Nuclear Physics B \textbf{190} (1981), 288--300,
\href{https://doi.org/10.1016/0550-3213(81)90560-5}{doi:10.1016/0550-3213(81)90560-5}.

\bibitem{Balaban1984}
T.~Balaban,
\emph{Propagators and renormalization transformations for lattice gauge
theories. I}, Communications in Mathematical Physics \textbf{95} (1984),
17--40,
\href{https://doi.org/10.1007/BF01215753}{doi:10.1007/BF01215753}.

\bibitem{Balaban1987}
T.~Balaban,
\emph{Renormalization group approach to lattice gauge field theories. I.
Generation of effective actions in a small field approximation and a coupling
constant renormalization in four dimensions}, Communications in Mathematical
Physics \textbf{109} (1987), 249--301,
\href{https://doi.org/10.1007/BF01215223}{doi:10.1007/BF01215223}.
\end{thebibliography}

\end{document}
